\documentclass[11pt]{article}
\usepackage[margin=1in]{geometry}
\usepackage[T1]{fontenc}
\usepackage{lmodern}
\usepackage{amsmath,amssymb,amsthm}
\usepackage{graphicx}
\usepackage{hyperref}

\newtheorem{theorem}{Theorem}
\newtheorem{lemma}{Lemma}
\newtheorem{remark}{Remark}
\newcommand{\M}{\mathbb M}
\newcommand{\C}{\mathbb C}
\newcommand{\Tr}{\operatorname{Tr}}
\newcommand{\rank}{\operatorname{rank}}

\title{The trace-arithmetic distance is a metric\\
on the full rank-one positive cone}
\author{Codex agent lane 6}
\date{August 9, 2026}

\begin{document}
\maketitle

\section*{Status and source question}

\textbf{Status: candidate partial result, likely valid.}

For $p>0$ and positive semidefinite matrices $A,B$, Zhang \cite{Zhang}
uses the operations
\[
 A\mathbin{\kappa_p}B
   =\left(A^{p/4}B^{p/2}A^{p/4}\right)^{1/p},
 \qquad
 d_p(A,B)^2
   =\Tr\left(\frac{A+B}{2}-A\mathbin{\kappa_p}B\right).
\]
Problem 1.2 on page 2 of the source asks whether $d_p$ is a metric on
$\M_n(\C)_+$ for every $1<p<2$, and asks the corresponding question for a
general $C^*$-algebra with a faithful trace.  Here is the source statement.

\begin{center}
  \includegraphics[width=0.98\linewidth]{figures/open_problem_crop.png}
\end{center}

The source gives a counterexample at $p=1/2$ but states that the full-cone
case $1<p<2$ remains open \cite[Remark 4.3]{Zhang}.  The result below proves
the metric property, throughout the entire requested parameter interval, on
the cone consisting of all positive multiples of rank-one projections, with
the zero matrix included.  In fact the proof works for every $0<p\leq2$.

\section*{Proof intuition}

Write a rank-one positive matrix as $aP$, where $a\geq0$ and $P$ is a
rank-one projection.  If $P=uu^*$ and $Q=vv^*$, the trace of the noncommutative
mean collapses to
\[
 \Tr(aP\mathbin{\kappa_p}bQ)
   =\sqrt{ab}\,|\langle u,v\rangle|^{2/p}.
\]
Thus the only directional datum is the Fubini--Study angle
$\theta(P,Q)=\arccos|\langle u,v\rangle|$.  For $\alpha=2/p\geq1$, deform
this angle by
\[
 f_\alpha(t)=\arccos((\cos t)^\alpha).
\]
The key one-variable fact is that $f_\alpha$ is increasing and concave on
$[0,\pi/2]$, hence it turns the projective angle into another metric.  After
putting $r=\sqrt a$ and $s=\sqrt b$, the distance becomes
\[
 \sqrt2\,d_p(aP,bQ)
   =\sqrt{r^2+s^2-2rs\cos(f_\alpha(\theta(P,Q)))}.
\]
This is exactly the Euclidean cone metric over the deformed projective metric.
The triangle inequality then follows by realizing any three base distances as
the angular distances of three unit vectors and applying the ordinary
Euclidean triangle inequality.  This argument controls the radial variables
$a,b$ that are absent from the previously known fixed-trace projector result.

The same radial decomposition also gives an exact reformulation of the full
problem.  On density matrices put
$F_p(\rho,\sigma)=\Tr(\rho\mathbin{\kappa_p}\sigma)$ and
$\delta_p=\arccos F_p$.  Then $d_p$ is a metric on the full positive cone if
and only if $\delta_p$ is a metric on density matrices.  Equivalently, every
three-state matrix $[F_p(\rho_i,\rho_j)]_{i,j=1}^3$ must be positive
semidefinite.  This isolates the unresolved part as one explicit trace
inequality; the packet does not prove that inequality, so the full cone remains
open.

\section*{Formal proof}

Let
\[
 \mathcal R_n^+
   :=\{0\}\cup\{aP:a>0,\ P\in\M_n(\C)_+,\ P^2=P,\ \rank P=1\}.
\]

\begin{lemma}[Projective-angle deformation]\label{lem:deformation}
Let $\theta(P,Q)=\arccos|\langle u,v\rangle|$ for $P=uu^*$ and $Q=vv^*$,
where $u,v$ are unit vectors.  If $\alpha\geq1$ and
\[
 f_\alpha(t)=\arccos((\cos t)^\alpha),\qquad 0\leq t\leq\frac\pi2,
\]
then
\[
 \delta_\alpha(P,Q):=f_\alpha(\theta(P,Q))
\]
is a metric on the rank-one projective space and takes values in
$[0,\pi/2]$.
\end{lemma}

\begin{proof}
The Fubini--Study angle $\theta$ is a metric.  For completeness, this follows
from the spherical triangle inequality: choose phases of successive unit
representatives so that the spherical angles of the first and second, and of
the second and third, are respectively $\theta(P,Q)$ and $\theta(Q,R)$.
The spherical angle between the first and third representatives is at least
the minimum over their relative phase, which is $\theta(P,R)$, and is at most
the sum of the other two angles.

It remains to record the properties of $f_\alpha$.  It is increasing,
$f_\alpha(0)=0$, and $f_\alpha(t)\geq t$, because
$(\cos t)^\alpha\leq\cos t$.  We now prove concavity.  On $(0,\pi/2)$,
\[
 f_\alpha'(t)
  =\frac{\alpha(\cos t)^{\alpha-1}\sin t}
          {\sqrt{1-(\cos t)^{2\alpha}}}.
\]
With $y=\cos^2t$, the square of $f_\alpha'(t)/\alpha$ is
\[
 G(y)=\frac{y^{\alpha-1}(1-y)}{1-y^\alpha}.
\]
After multiplication by the positive denominator, the sign of $G'(y)$ is the
sign of
\[
 h(y)=\alpha-1-\alpha y+y^\alpha.
\]
Since $h(1)=0$ and
$h'(y)=\alpha(y^{\alpha-1}-1)\leq0$ on $[0,1]$, one has $h(y)\geq0$.
Hence $G$ is increasing in $y$.  Since $y=\cos^2t$ decreases with $t$,
$f_\alpha'$ decreases, so $f_\alpha$ is concave on $[0,\pi/2]$.

Concavity and $f_\alpha(0)=0$ imply subadditivity whenever
$x,y,x+y\in[0,\pi/2]$: indeed,
\[
 f_\alpha(x)\geq\frac{x}{x+y}f_\alpha(x+y),\qquad
 f_\alpha(y)\geq\frac{y}{x+y}f_\alpha(x+y).
\]
For three projections put $x=\theta(P,Q)$ and $y=\theta(Q,R)$.  If
$x+y\leq\pi/2$, monotonicity, the triangle inequality for $\theta$, and
subadditivity give
\[
 \delta_\alpha(P,R)\leq f_\alpha(x+y)
 \leq f_\alpha(x)+f_\alpha(y).
\]
If $x+y>\pi/2$, then
\[
 \delta_\alpha(P,R)\leq\frac\pi2<x+y
 \leq f_\alpha(x)+f_\alpha(y).
\]
Definiteness and symmetry are immediate.  Thus $\delta_\alpha$ is a metric.
\end{proof}

\begin{lemma}[Elementary metric-cone lemma]\label{lem:cone}
Let $(X,\delta)$ be a metric space with $0\leq\delta\leq\pi/2$.  Identify all
points $(0,x)$ in $[0,\infty)\times X$, and define
\[
 D((r,x),(s,y))^2=r^2+s^2-2rs\cos\delta(x,y).
\]
Then $D$ is a metric on the resulting cone.
\end{lemma}

\begin{proof}
Only the triangle inequality needs proof.  Given $x,y,z\in X$, put
\[
 a=\delta(x,y),\qquad b=\delta(y,z),\qquad c=\delta(x,z).
\]
The metric inequalities give $|a-b|\leq c\leq a+b$, and $a+b\leq\pi$.
Consequently there are unit vectors $U,V,W\in\mathbb R^3$ whose pairwise
angles are $a,b,c$.  Explicitly, when $\sin a\sin b\ne0$, take
\[
 V=(1,0,0),\quad U=(\cos a,\sin a,0),
\]
\[
 W=(\cos b,\sin b\cos\gamma,\sin b\sin\gamma),\qquad
 \cos\gamma=\frac{\cos c-\cos a\cos b}{\sin a\sin b}.
\]
The inequalities $|a-b|\leq c\leq a+b$ say exactly that
$|\cos\gamma|\leq1$; the cases with a zero sine follow by a limit or directly.
For nonnegative $r,s,t$, the Euclidean triangle inequality now gives
\[
 \|rU-tW\|_2\leq\|rU-sV\|_2+\|sV-tW\|_2,
\]
which is precisely
\[
 D((r,x),(t,z))\leq D((r,x),(s,y))+D((s,y),(t,z)).
\]
Symmetry and nonnegativity are clear.  Equality $D=0$ identifies exactly the
same positive-radius base point, together with the already identified
zero-radius points.
\end{proof}

\begin{theorem}[Exact reduction of the full-cone problem]\label{thm:reduction}
Fix $n\geq2$ and a value of $p$ for which $d_p$ is a definite nonnegative
function.  Let
\[
 \mathcal D_n=\{\rho\in\M_n(\C)_+:\Tr\rho=1\},\qquad
 F_p(\rho,\sigma)=\Tr(\rho\mathbin{\kappa_p}\sigma),
\]
and set $\delta_p(\rho,\sigma)=\arccos F_p(\rho,\sigma)$.  The following are
equivalent:
\begin{enumerate}
  \item $d_p$ is a metric on $\M_n(\C)_+$;
  \item $\delta_p$ is a metric on $\mathcal D_n$;
  \item for all $\rho,\sigma,\omega\in\mathcal D_n$,
  \begin{equation}\label{eq:gram}
  1+2F_p(\rho,\sigma)F_p(\sigma,\omega)F_p(\omega,\rho)
  \geq F_p(\rho,\sigma)^2+F_p(\sigma,\omega)^2
       +F_p(\omega,\rho)^2.
  \end{equation}
\end{enumerate}
For $1<p<2$, therefore, the matrix part of Problem 1.2 is exactly the
three-state trace inequality \eqref{eq:gram}.
\end{theorem}

\begin{proof}
The scalar overlap is symmetric because
\[
 F_p(\rho,\sigma)
 =\Tr\bigl|\sigma^{p/4}\rho^{p/4}\bigr|^{2/p},
\]
and an operator and its adjoint have the same singular values.  The assumed
nonnegativity and definiteness of $d_p$ give
$0\leq F_p(\rho,\sigma)\leq1$, with equality at the upper endpoint exactly
when $\rho=\sigma$.  Thus $\delta_p$ is a well-defined symmetric definite
function with values in $[0,\pi/2]$.

For nonzero $A,B$, write
\[
 A=a\rho,\qquad B=b\sigma,\qquad
 a=\Tr A,\quad b=\Tr B,\quad \rho,\sigma\in\mathcal D_n.
\]
Separate homogeneity of the mean gives
\[
 \Tr(A\mathbin{\kappa_p}B)=\sqrt{ab}\,F_p(\rho,\sigma),
\]
and hence
\begin{equation}\label{eq:cone-decomposition}
 2d_p(A,B)^2
 =a+b-2\sqrt{ab}\cos\delta_p(\rho,\sigma).
\end{equation}
Thus $\sqrt2d_p$ is precisely the cosine-cone distance with radial coordinate
$\sqrt{\Tr A}$.  Lemma~\ref{lem:cone} proves that (2) implies (1).

Conversely, suppose that the triangle inequality for $\delta_p$ fails.  After
relabeling, there are density matrices $\rho,\sigma,\omega$ such that
\[
 \gamma:=\delta_p(\rho,\omega)
 >\alpha+\beta,\qquad
 \alpha:=\delta_p(\rho,\sigma),\quad
 \beta:=\delta_p(\sigma,\omega).
\]
All these angles lie in $[0,\pi/2]$, and strict failure forces
$\alpha,\beta>0$.  In the Euclidean plane let $X,Y,Z$ be unit vectors on rays
of angles $0,\alpha,\alpha+\beta$, respectively.  The line segment from $X$
to $Z$ meets the intermediate ray at $sY$ for some $s>0$.  Apply
\eqref{eq:cone-decomposition} to
\[
 A=\rho,\qquad B=s^2\sigma,\qquad C=\omega.
\]
Because $sY$ lies on the segment $[X,Z]$,
\begin{align*}
 \sqrt2\bigl(d_p(A,B)+d_p(B,C)\bigr)
 &=\|X-sY\|_2+\|sY-Z\|_2\\
 &=\|X-Z\|_2
 =\sqrt{2-2\cos(\alpha+\beta)}\\
 &<\sqrt{2-2\cos\gamma}
 =\sqrt2d_p(A,C).
\end{align*}
Thus failure of the angular triangle inequality forces failure on the positive
cone, proving (1) implies (2).

It remains to identify (2) with (3).  Put
\[
 x=F_p(\rho,\sigma)=\cos\alpha,\quad
 y=F_p(\sigma,\omega)=\cos\beta,\quad
 z=F_p(\omega,\rho)=\cos\gamma.
\]
The left side minus the right side of \eqref{eq:gram} is the determinant of
the symmetric matrix with diagonal entries $1$ and off-diagonal entries
$x,y,z$.  Moreover,
\[
 1+2xyz-x^2-y^2-z^2
 =\sin^2\alpha\sin^2\beta
  -(\cos\gamma-\cos\alpha\cos\beta)^2.
\]
Since $\alpha,\beta,\gamma\in[0,\pi/2]$, nonnegativity is equivalent to
$|\alpha-\beta|\leq\gamma\leq\alpha+\beta$, exactly the three triangle
inequalities.  This proves the remaining equivalence.
\end{proof}

\begin{theorem}[Metricity on the full rank-one positive cone]
For every $n\geq2$ and every $0<p\leq2$, the restriction of $d_p$ to
$\mathcal R_n^+$ is a metric.  In particular, this holds for every
$1<p<2$, the parameter range left open on the full cone in Problem 1.2 of
\cite{Zhang}.
\end{theorem}

\begin{proof}
Let $A=aP$ and $B=bQ$, with $a,b>0$, $P=uu^*$, $Q=vv^*$, and $u,v$ unit
vectors.  Functional calculus and $PQP=|\langle u,v\rangle|^2P$ yield
\[
 A^{p/4}B^{p/2}A^{p/4}
  =(ab)^{p/2}|\langle u,v\rangle|^2P.
\]
Therefore
\[
 \Tr(A\mathbin{\kappa_p}B)
  =\sqrt{ab}\,|\langle u,v\rangle|^{2/p}.
\]
The same formula, by continuity, also covers either $a=0$ or $b=0$.
Put
\[
 r=\sqrt a,\qquad s=\sqrt b,\qquad
 \alpha=\frac2p\geq1,
\]
and let $\theta(P,Q)=\arccos|\langle u,v\rangle|$.  Then
\begin{align*}
 2d_p(aP,bQ)^2
 &=r^2+s^2-2rs(\cos\theta(P,Q))^\alpha\\
 &=r^2+s^2-2rs\cos\delta_\alpha(P,Q).
\end{align*}
By Lemma~\ref{lem:deformation}, $\delta_\alpha$ is a metric with diameter at
most $\pi/2$.  Lemma~\ref{lem:cone} therefore says that the square root of the
last expression is a metric on the cone.  Since this square root equals
$\sqrt2\,d_p$, the same is true of $d_p$.  The cone point at radius zero is
exactly the zero matrix, so definiteness has the intended matrix meaning.
\end{proof}

\section*{Verification, scope, and novelty check}

The proofs are analytic; the script \texttt{code/check\_rank\_one\_metric.py}
only checks the trace formula and random triangle instances and is not part of
the proof.  The crucial reviewer checks are the sign calculation proving
concavity of $f_\alpha$, the use of the elementary metric-cone lemma, and both
directions of Theorem~\ref{thm:reduction}.

This result does \emph{not} settle Problem 1.2 on the full matrix cone or for a
general $C^*$-algebra.  The noncommutative difficulty begins once matrices have
rank at least two.  Theorem~\ref{thm:reduction} isolates that difficulty as the
three-state Gram determinant inequality \eqref{eq:gram}, which is not proved
here.

A bounded novelty search was performed before promotion:
\begin{itemize}
  \item the local run indexes and parsed source collection were searched by
  arXiv id; related searches used rank-one matrices, projectors, trace
  arithmetic, and the symmetric Bures--Wasserstein divergence;
  \item Kom\'alovics--Moln\'ar \cite{KM} prove the threshold only on the
  normalized rank-one projector set;
  \item a targeted web/full-text check of Vuong's April 2026 follow-up
  \cite{Vuong} found the commuting cone and the normalized rank-one projector
  set, and explicitly says that full-cone metricity for the corresponding
  parameters remains open;
  \item no explicit theorem for arbitrary scalar multiples $aP$ was found.
\end{itemize}
The extension may be regarded as a natural metric-cone completion of the known
projector theorem, so novelty confidence is moderate rather than high.

\textbf{Human-review recommendation.}  Promote after checking
Lemma~\ref{lem:deformation}, especially the monotonicity direction in the
$G(y)$ calculation, and confirming that the searched literature does not use
``rank-one cone'' to include arbitrary non-normalized scalar multiples.

\begin{thebibliography}{9}

\bibitem{Zhang}
T. Zhang,
\emph{Trace arithmetic--$\kappa_p$ inequality},
arXiv:2602.11922, 2026.

\bibitem{KM}
\'A. Kom\'alovics and L. Moln\'ar,
\emph{On a parametric family of distance measures that includes the Hellinger
and the Bures distances},
J. Math. Anal. Appl. 529 (2024), no. 2, article 127226.

\bibitem{Vuong}
T.-D. Vuong,
\emph{Quantitative geometry of the symmetric $\alpha$-$z$
Bures--Wasserstein divergence},
Filomat, prepublication manuscript, April 2026.

\end{thebibliography}

\end{document}
